\documentclass[10pt,a4paper]{amsart}
\usepackage[margin=19mm]{geometry}
\usepackage{amsmath,amssymb,mathtools}
\usepackage[scr=rsfs,cal=euler]{mathalfa}
\usepackage{xfrac}
\usepackage[unicode,hidelinks,bookmarks=false]{hyperref}
\newcommand{\ie}{i.e.\ }
\newcommand{\cf}{cf.\ }
\newcommand{\resp}{respectively\ }
\newcommand{\Subsection}{\subsection}
\providecommand{\nobreakdash}{\nobreak\hbox{-}}
\setlength{\emergencystretch}{2em}
\multlinegap0pt
\usepackage{amsthm}
\usepackage{xcolor}
\providecommand{\newblock}{\space}
\newtheorem{thm}{Theorem}[section]
\newtheorem{lm}[thm]{Lemma}
\newtheorem{defn}[thm]{Definition}
\newtheorem{notation}[thm]{Notation}
\newtheorem{prop}[thm]{Proposition}
\newtheorem{conj}[thm]{Conjecture}
\newtheorem{cor}[thm]{Corollary}
\newtheorem{claim}{Claim}
\newtheorem{pf}{Proof}
\newtheorem{rmk}[thm]{Remark}
\newcommand{\bigO}{\mathcal{O}}
\newcommand{\Broots}{\mathcal{R}}
\newcommand{\Cb}[1]{\textbf{\textcolor{blue}{{#1}}}}
\newcommand{\CC}{\mathcal{C}}
\newcommand{\Ch}{\mathcal{C}}
\newcommand{\complexC}{\mathbb{C}}
\newcommand{\conf}{\mathcal{X}}
\newcommand{\Cr}[1]{\textbf{\textcolor{red}{{#1}}}}

\newcommand{\ddbar}[2]{\frac{{\mathrm d}#1}{2\pi {\mathrm i}#2}}
\newcommand{\dg}{\mathrm{dist}}
\newcommand{\dtt}{\text{dist}}
\newcommand{\eigl}{\mathcal{L}}
\newcommand{\eigr}{\mathcal{R}}
\newcommand{\exe}{\mathrm{err}}
\newcommand{\gd}{{\mathcal{G}}}
\newcommand{\ii}{\mathrm{i}}
\newcommand{\inn}{\mathrm{in}}
\newcommand{\intZ}{\mathbb{Z}}
\newcommand{\limp}{\mathrm{p}}
\newcommand{\LL}{\mathrm{L}}
\newcommand{\out}{\mathrm{out}}
\newcommand{\prob}{\mathbb{P}}
\newcommand{\realR}{\mathbb{R}}
\newcommand{\rg}{\mathrm{g}}
\newcommand{\RR}{\mathrm{R}}
\newcommand{\rr}{\mathrm{r_{max}}}
\newcommand{\rz}{\mathrm{z}}
\newcommand{\wt}{\mathcal{Q}}
\title[One-point distribution of the geodesic in directed last passage percolation]{One-point distribution of the geodesic in directed last passage percolation: the generalized Cauchy identity (Lemma 2.9)}
\author{Zhipeng Liu}
\date{Source: arXiv:2105.15150v3 (CC BY 4.0)}
\begin{document}
\maketitle
\noindent\emph{Source and licence.} One-point distribution of the geodesic in directed last passage percolation, arXiv:2105.15150v3, distributed under the Creative Commons Attribution 4.0 International licence, http://creativecommons.org/licenses/by/4.0/, as recorded in the arXiv OAI metadata for this exact version and shown on the abstract page https://arxiv.org/abs/2105.15150v3 (v3 of 31 August 2025). Full rights notice: licenses/arxiv-2105.15150-CC-BY-4.0.txt.\par\smallskip
\noindent\textit{Editorial scope.} Counted unit: the generalized Cauchy identity, the author\textquotesingle s own lemma with source label \textbackslash{}label\{lem:Cauchy\_generalization\} (source0.tex lines 1832--1847), printed as Lemma 2.9 of the article, delivered with its complete original proof, the whole of the author\textquotesingle s \textbackslash{}begin\{proof\} block at source0.tex lines 1849--2050, as the single primary proof body. Nothing is abridged: the proof body is delivered from its opening \textbackslash{}begin\{proof\}[Proof of Lemma\textasciitilde{}\textbackslash{}ref\{lem:Cauchy\_generalization\}] line to its closing \textbackslash{}end\{proof\}, including the auxiliary sums $C_{p,q}$, the residue computation of $C_{-1,2}$, Table 1 with its six tabulated $C_{p,q}$ values, and the final simplification. Required context retained uncounted: source0.tex lines 941--950, the display that defines $\hat H$ (source label eq:hat\_H), which the counted statement names in its last line; it is the only label that the statement and proof reference from outside, and its own text references nothing further.\par
\noindent\textit{Wrapper and numbering.} The article\textquotesingle s \textbackslash{}documentclass\{article\} and its package list are replaced by the AMS amsart wrapper of the pinned TeX Live runtime; the author\textquotesingle s own theorem environments (thm, lm, defn, notation, prop, conj, cor, claim, pf, rmk, source lines 198--207) and his own macros (source lines 210--238) are carried over verbatim, and one \textbackslash{}providecommand shim supplies \textbackslash{}newblock, which the AMS wrapper does not itself define. The tikz/pgfplots/drawing block of the author\textquotesingle s preamble is not carried over: no drawing, figure or graphical construct lies in the delivered spans. Two editorial counter commands immediately before the counted statement (\textbackslash{}setcounter\{section\}\{2\} and \textbackslash{}setcounter\{thm\}\{8\}) restore the article\textquotesingle s own number, so the statement prints as Lemma 2.9 exactly as in the version PDF (page 20). The author\textquotesingle s \textbackslash{}numberwithin\{equation\}\{section\} is not carried over, so the displays of this unit carry the wrapper\textquotesingle s own consecutive numbers rather than the article\textquotesingle s section-based equation numbers; every internal cross-reference (the lemma itself, its displayed identity, the five auxiliary displays, Table 1 and the two definitions) resolves to the correct object inside the unit.\par
\noindent\textit{Citation map.} The delivered unit invokes exactly one citation, \textbackslash{}cite[Lemma 5.5]\{Baik-Liu17\} at source line 1881. Its entry is transcribed from the article\textquotesingle s own in-source thebibliography (lines 3930--3933) and prints with the article\textquotesingle s own label [BL19]; the delivered bibliography contains that entry and no other, which is exactly the set of keys the delivered text cites.\par
\noindent\textit{Assets.} The e-print of this version ships one file besides the author\textquotesingle s .tex: contour\_figure.pdf, the illustration of Figure 1 of the article. No retained span contains any \textbackslash{}includegraphics, \textbackslash{}input or \textbackslash{}include command, the delivered TeX is a single self-contained file, and no figure, artwork, class, style, font or bitmap is redistributed in the package. Every mathematical symbol, sentence, table entry and label of the delivered spans is the author\textquotesingle s own native text; no mathematics was written, repaired, re-derived or normalized, and no printed slip inside these spans was corrected.\par
\begin{equation}
	\label{eq:hat_H}
	\hat H\left(W;W'\right):= 
	   \frac{1}{2} \left(
	                   \sum_i w_i -\sum_{i'} w'_{i'}  
	               \right)^2
	   -\frac{1}{2} \left(
	                  \sum_i w_i^2 -\sum_{i'} \left(w'_{i'}\right)^2  
	                \right)
\end{equation}

\setcounter{section}{2}
\setcounter{thm}{8}
\begin{lm}
	\label{lem:Cauchy_generalization}
	Suppose $X=(x_1,\cdots,x_N)$ and $Y=(y_1,\cdots,y_N)$ are two vectors in $\complexC^N$ satisfying $x_i\ne y_j$ for all $1\le i,j\le N$. Suppose $\rz$ is an arbitrary complex number. Then we have the following identity
	\begin{equation}
		\label{eq:Cauchy_generalization}
		\begin{split}
		&\sum_{a,b=1}^N (-1)^{a +b}
		   y_b \det\left[ \frac{\rz}{ x_i-y_j} \frac{y_j}{x_i}
		                        +\frac{1}{-x_i+y_j}\frac{x_i}{y_j}   
		                  \right]_{\substack{i\ne a\\ j\ne b}}\\
		                  &
		=(1-\rz)^{N-2}\left(\hat H(X;Y)+\rz\prod_{i=1}^N\frac{y_i}{x_i}\hat H(Y;X)\right)\det\left[\frac{1}{y_j-x_i}\right]_{i,j=1}^N,
		\end{split}
	\end{equation}
where $\hat H$ is defined in~\eqref{eq:hat_H}.
\end{lm}

\begin{proof}[Proof of Lemma~\ref{lem:Cauchy_generalization}]
	We first use the identity
	\begin{equation*}
		\frac{\rz}{ x-y} \frac{y}{x}
		+\frac{1}{-x+y}\frac{x}{y} 
		=(1-\rz)\frac{x}{y}
		    \cdot\left(\frac{1}{-x+y} 
		              -\frac{\rz}{1-\rz}\frac{1}{x}
		              -\frac{\rz}{1-\rz}\frac{y}{x^2}
		         \right)
	\end{equation*}
and write the left hand side of~\eqref{eq:Cauchy_generalization} as
\begin{equation}
	(1-\rz)^{N-1}\prod_{i=1}^N\frac{x_i}{y_i} \sum_{a,b=1}^N(-1)^{a+b}\frac{y_b^2}{x_a}
	\det\left[\frac{1}{-x_i+y_j} 
	-\frac{\rz}{1-\rz}\frac{1}{x_i}
	-\frac{\rz}{1-\rz}\frac{y_j}{x_i^2}\right]_{\substack{i\ne a \\ j\ne b}}.
\end{equation}
Thus the equation~\eqref{eq:Cauchy_generalization} is equivalent to, by setting $u=-\rz/(1-\rz)$,
\begin{equation}
	\label{eq:sum_identity_01}
	\begin{split}
	&\sum_{a,b=1}^N (-1)^{a+b} \frac{y_b^2}{x_a} 
	       \det\left[  \frac{1}{-x_i+y_j} +u \frac{1}{x_i} +u \frac{y_j}{x_i^2} \right]_{\substack{i\ne a\\ j\ne b}}\\
	       &
	       = -\left(
	        (u-1)\prod_{i=1}^N \frac{y_i}{x_i} \hat H(X;Y)
	        + u \prod_{i=1}^N\frac{y_i^2}{x_i^2}\hat H(Y; X)
	       \right)\cdot \det\left[\frac{1}{y_j-x_i}\right]_{i,j=1}^N.
	       \end{split}
\end{equation}

The proof of~\eqref{eq:sum_identity_01} is tedious while the strategy is quite straightforward. Below we will show the proof but omit some details which are direct to check. We remark that the strategy was applied to a much simpler identity in \cite[Lemma 5.5]{Baik-Liu17}, but this identity~\eqref{eq:sum_identity_01} is much more complicated.

Before we prove~\eqref{eq:sum_identity_01}, we need to prepare some easier identities. We denote 
\begin{equation*}
	X(w):= \prod_{i=1}^N (w-x_i),\quad Y(w):= \prod_{i=1}^N (w-y_i),
\end{equation*}
and introduce
\begin{equation*}
	C_{p,q} = \sum_{a,b=1}^n x_a^py_b^q\frac{Y(x_a)X(y_b)}{(x_a-y_b)X'(x_a)Y'(y_b)},
\end{equation*}
where $p,q$ are both integers. It is not hard to verify, by using the Cauchy determinant formula, that
\begin{equation}
	\label{eq:C_ab_relation_Cauchy}
	C_{p,q} = \sum_{a,b=1}^N (-1)^{a+b} x_a^p y_b^q \det\left[\frac{1}{-x_i+y_j}\right]_{\substack{i\ne a\\ j\ne b}}
	/\det\left[\frac{1}{-x_i+y_j}\right]_{i,j=1}^N.
\end{equation}

One can evaluate $C_{p,q}$ by converting the sum as a residue computation of an integral on the complex plane. As an illustration, we show how to obtain $C_{-1,2}$, then we will list all the $C_{p,q}$ values we will use later without providing proofs, see Table~\ref{table:C}.

We consider a double integral
\begin{equation*}
	\int_{|y|=R_2} \int_{|x|=R_1} \frac{y^2}{x} \frac{Y(x)X(y)}{(x-y)X(x)Y(y)} \ddbar{x}{} \ddbar{y}{},
\end{equation*}
where $R_1>R_2> \max_i\{|x_i|+|y_i|\}$. Note that we can deform the $x$-contour to infinity and the integral becomes zero. Hence the above double integral is zero. On the other hand, we can change the order of integrals and evaluate the $y$-integral first. It gives a sum over all roots of $Y(y)$:
\begin{equation*}
	0= \int_{|x|=R_1} \sum_{b=1}^N \frac{y_b^2}{x} \frac{Y(x)X(y_b)}{(x-y_b) X(x) Y'(y_b)} \ddbar{x}{}.
\end{equation*}
Then we exchange the summation and integral, and evaluate the $x$-integral by computing the residues within the contour. Note that $x=y_b$ is not a pole. We get
\begin{equation}
	\label{eq:aux_01}
	0=C_{-1,2} -\frac{Y(0)}{X(0)} \sum_{b=1}^N y_b\frac{X(y_b)}{Y'(y_b)}.
\end{equation}
We need to continue to evaluate the summation in~\eqref{eq:aux_01}. We have, by a residue computation,
\begin{equation*}
	\begin{split}
	\sum_{b=1}^N y_b\frac{X(y_b)}{Y'(y_b)}&= \int_{|y|=R_2}y\frac{X(y)}{Y(y)}\ddbar{y}{}\\
	&= \int_{|y|=R_2} y\left(1+\frac{1}{y}\sum_{i=1}^N(y_i-x_i)+\frac{1}{y^2}\hat H(X;Y) +O(y^{-3})\right)\ddbar{y}{}
	 = \hat H(X;Y),
\end{split}
\end{equation*}
where we evaluated the integral by expanding the integrand for large $y$. Here the function $\hat H$ is defined in~\eqref{eq:hat_H}. 

By inserting the above formula to~\eqref{eq:aux_01}, we obtain
\begin{equation*}
	C_{-1,2}=\prod_i\frac{y_i}{x_i}\hat H(X;Y).
\end{equation*}
Using similar calculations, we can find all $C_{p,q}$ for small $p,q$ values. In Table~\ref{table:C} we list some $C_{p,q}$ identities we will use in the proof of~\eqref{eq:sum_identity_01}. We remark that the proof of these identities are analogous to that of $C_{-1,2}$ without adding extra difficulties.

\begin{table}
\begin{tabular}{ |p{1.7cm}|p{2.5cm}|p{1.8cm}|p{7.5cm}|  }
    \hline
    \rule{0pt}{0.6cm}
	Expression& Value &Expression&Value\\
	&&&\\
	\hline
	\rule{0pt}{0.6cm}$C_{0,-1}$   & $\displaystyle 1-\prod_i \frac{x_i}{y_i}$    &
	$C_{-1,2}$ &$\displaystyle \prod_i\frac{y_i}{x_i} \hat H(X;Y)$
	\\
    \hline
	\rule{0pt}{0.6cm}
	$C_{-1,0}$& $\displaystyle-1+ \prod_i \frac{y_i}{x_i}$&$C_{-1,1}-C_{0,0}$ & $\displaystyle \left(1-\prod_i\frac{y_i}{x_i}\right)\sum_i(x_i-y_i)$
	\\
	\hline
		\rule{0pt}{0.6cm}
	$C_{1,0}$ & $\displaystyle -\hat H(Y;X)$
	& 	$C_{0,2}-C_{1,1}$
	&$\displaystyle -\sum_i(x_i-y_i) \hat H(X;Y)$
	\\
	\hline
			\rule{0pt}{0.6cm}
			$C_{0,1}$& $\displaystyle \hat H(X;Y)$&$C_{-2,1}$& $\displaystyle -1 + \prod_i\frac{y_i}{x_i}\left( 1- \sum_i\left(\frac{1}{x_i}-\frac{1}{y_i}\right)\sum_i\left(x_i-y_i\right)\right)$\\
	\hline
\end{tabular}
\caption{\label{table:C}Values of some $C_{p,q}$ expressions.}
\end{table}


We need to evaluate
\begin{equation*}
	\det\left[\frac{1}{-x_i+y_j} +u\frac{1}{x_i}\right]_{i,j=1}^N =\det\left[\frac{1}{-x_i+y_j}\right]_{i,j=1}^N + u\sum_{a,b=1}^N (-1)^{a+b}\frac{1}{x_a} \det\left[\frac{1}{-x_i+y_j}\right]_{\substack{i\ne a\\ j\ne b}}.
\end{equation*}
By applying~\eqref{eq:C_ab_relation_Cauchy} and finding the $C_{-1,0}$ value in Table~\ref{table:C}, we get
\begin{equation}
	\label{eq:aux_02}
	\begin{split}
	\det\left[\frac{1}{-x_i+y_j} +u\frac{1}{x_i}\right]_{i,j=1}^N 
	&=\det\left[\frac{1}{-x_i+y_j}\right]_{i,j=1}^N \left(1+uC_{-1,0}\right)\\
	&=\det\left[\frac{1}{-x_i+y_j}\right]_{i,j=1}^N  \left(1+ u\left(-1+\prod_{i=1}^N\frac{y_i}{x_i}\right)\right).
	\end{split}
\end{equation}
Then we evaluate
\begin{equation*}
	\det\left[  \frac{1}{-x_i+y_j} +u \frac{1}{x_i} +u \frac{y_j}{x_i^2} \right]_{i,j=1}^N = \det\left[\frac{1}{-x_i+y_j} +u\frac{1}{x_i}\right]_{i,j=1}^N
	+u\sum_{a,b=1}^N (-1)^{a+b}\frac{y_b}{x_a^2}\det\left[\frac{1}{-x_i+y_j} +u\frac{1}{x_i}\right]_{\substack{i\ne a\\ j\ne b}}.
\end{equation*}
We insert~\eqref{eq:aux_02} in the above equation and obtain
\begin{equation}
	\label{eq:aux_03}
	\begin{split}
		&\det\left[  \frac{1}{-x_i+y_j} +u \frac{1}{x_i} +u \frac{y_j}{x_i^2} \right]_{i,j=1}^N\\
		&= \det\left[\frac{1}{-x_i+y_j}\right]_{i,j=1}^N  \left(1- u+ u\prod_{i=1}^N\frac{y_i}{x_i}\right)
		+u(1-u)\sum_{a,b=1}^N (-1)^{a+b}\frac{y_b}{x_a^2}\det\left[\frac{1}{-x_i+y_j}\right]_{\substack{i\ne a\\ j\ne b}}\\
		&\quad +u^2\prod_{i=1}^N\frac{y_i}{x_i}\sum_{a,b=1}^N (-1)^{a+b}\frac{1}{x_a}\det\left[\frac{1}{-x_i+y_j}\right]_{\substack{i\ne a\\ j\ne b}}\\
		&=\det\left[\frac{1}{-x_i+y_j}\right]_{i,j=1}^N  \left(1- u+ u\prod_{i=1}^N\frac{y_i}{x_i}
		+u(1-u) C_{-2,1} + u^2\prod_{i=1}^N\frac{y_i}{x_i} C_{-1,0}\right).
	\end{split}
\end{equation}
By inserting the values of $C_{-2,1}$ and $C_{-1,0}$ and simplifying the expression, we obtain
\begin{equation}
	\label{eq:aux_04}
	\begin{split}
		&\det\left[  \frac{1}{-x_i+y_j} +u \frac{1}{x_i} +u \frac{y_j}{x_i^2} \right]_{i,j=1}^N =  \det\left[\frac{1}{-x_i+y_j}\right]_{i,j=1}^N\\
		&\quad \cdot \left[
		1 - 2u\left( 1-\prod_{i=1}^N \frac{y_i}{x_i}\right) -(u-u^2)\prod_{i=1}^N \frac{y_i}{x_i}\sum_{i=1}^N\left(\frac{1}{x_i}-\frac{1}{y_i}\right)\sum_{i=1}^N (x_i-y_i) + u^2 \left(\prod_{i=1}^N \frac{y_i}{x_i}-1\right)^2
		\right].
	\end{split}
\end{equation}
Finally we are ready to prove~\eqref{eq:sum_identity_01}. Inserting~\eqref{eq:aux_04}, we can write
\begin{equation*}
	\begin{split}
		&\sum_{a,b=1}^N (-1)^{a+b} \frac{y_b^2}{x_a} 
		\det\left[  \frac{1}{-x_i+y_j} +u \frac{1}{x_i} +u \frac{y_j}{x_i^2} \right]_{\substack{i\ne a\\ j\ne b}}\\
		&=\sum_{a,b=1}^N (-1)^{a+b} \frac{y_b^2}{x_a} \det\left[\frac{1}{-x_i+y_j}\right]_{\substack{i\ne a\\ j\ne b}}
		\cdot \left[
		1 - 2u\left( 1-\frac{x_a}{y_b}\prod_{i=1}^N \frac{y_i}{x_i}\right)
		 + u^2 \left(\frac{x_a}{y_b}\prod_{i=1}^N \frac{y_i}{x_i}-1\right)^2\right.\\
		&\quad\left. -(u-u^2)\frac{x_a}{y_b}\prod_{i=1}^N \frac{y_i}{x_i}
		\left(-\frac{1}{x_a}+\frac{1}{y_b}+
		\sum_{i=1}^N\left(\frac{1}{x_i}-\frac{1}{y_i}\right)\right)
		\left(-x_a+y_b+ \sum_{i=1}^N (x_i-y_i)\right)
		\right].
	\end{split}
\end{equation*}
We apply~\eqref{eq:C_ab_relation_Cauchy} and rewrite the above equation as
\begin{equation}
	\begin{split}
		&\sum_{a,b=1}^N 
		    (-1)^{a+b} \frac{y_b^2}{x_a} 
		        \det\left[
		               \frac{1}{-x_i+y_j} +u \frac{1}{x_i} +u \frac{y_j}{x_i^2} 
		            \right]_{\substack{i\ne a\\ j\ne b}}
		       /\det\left[
		              \frac{1}{-x_i+y_j} +u \frac{1}{x_i} +u \frac{y_j}{x_i^2} 
		            \right]_{i,j=1}^N\\
	  &=\left(
	         (1-u)^2 +u(1-u)\prod_i\frac{y_i}{x_i}
	    \right)C_{-1,2}
	    +u\prod_i\frac{y_i}{x_i}
	       \left(1-u+u\cdot\prod_i\frac{y_i}{x_i}\right)
	        C_{1,0}\\
	  &\quad-
		  u(1-u)\prod_i\frac{y_i}{x_i}
		      \left(\sum_i\frac{1}{x_i}-\sum_i\frac{1}{y_i}\right)
		      \left(\sum_ix_i-\sum_iy_i\right) C_{0,1}\\
	  &\quad -
	     u(1-u)\prod_i\frac{y_i}{x_i}
	         \left(\sum_ix_i-\sum_iy_i\right)C_{0,0}
	        +
	     u(1-u)\prod_i\frac{y_i}{x_i}
	         \left(\sum_ix_i-\sum_iy_i\right)C_{-1,1}\\
	  &\quad -
	    u(1-u)\prod_i\frac{y_i}{x_i}
	         \left(\sum_i\frac{1}{x_i}-\sum_i\frac{1}{y_i}\right)C_{0,2}
	         +
	    u(1-u)\prod_i\frac{y_i}{x_i}
	         \left(\sum_i\frac{1}{x_i}-\sum_i\frac{1}{y_i}\right)C_{1,1}.
	\end{split}
\end{equation}
By checking the values of Table~\ref{table:C}, and noting that $\left(\sum_i (x_i-y_i)\right)^2=\hat H(X;Y)+\hat H(Y;X)$, we can simplify the above expression. It turns out, after a careful but straightforward calculation, the $u^2$ term vanishes, and the remaining terms match the right hand side of~\eqref{eq:sum_identity_01}. We hence complete the proof.
\end{proof}

\begin{thebibliography}{99}
\bibitem[BL19]{Baik-Liu17} Jinho Baik and Zhipeng Liu. \newblock Multipoint distribution of periodic {TASEP}. \newblock {\em J. Amer. Math. Soc.}, 32(3):609--674, 2019.
\end{thebibliography}
\end{document}
